\documentclass[11pt]{article}
\usepackage[a4paper,margin=1in]{geometry}
\usepackage{amsmath,amssymb,hyperref}
\title{An integer-matrix separation proving Conjecture 00000006330}
\author{gaochengzhecpu}
\date{October 3, 2026}
\begin{document}
\maketitle

The source asks for two matrices with identical determinant and permanent
but different integer types, with the separation described through
unimodular row transformations. The following pair works even if both
unimodular row and column operations are allowed:
\[
 A=\begin{pmatrix}1&0\\0&4\end{pmatrix},\qquad
 B=\begin{pmatrix}2&0\\0&2\end{pmatrix}.
\]
They satisfy
\[
 \det A=\det B=4,\qquad
 \operatorname{per}A=\operatorname{per}B=4.
\]
Both matrices are already in Smith normal form: their diagonal entries
are positive and the first divides the second. Their invariant-factor
lists are respectively $(1,4)$ and $(2,2)$.

A direct argument establishes the distinction without invoking the Smith
normal form classification theorem. For arbitrary integer matrices $U,V$,
every entry of $UBV=2UV$ is even. Therefore $UBV\ne A$, whose top-left
entry is one. In particular, there are no unimodular integer $U,V$ with
$UBV=A$. Thus the matrices have distinct integer-equivalence types.

The row-only interpretation is also covered in both orientations.
An equation $UB=A$ is impossible by parity. An equation $UA=B$ would
require $4U_{22}=2$ by comparing the bottom-right entries, which is
impossible over the integers. Hence neither matrix is obtained from the
other by unimodular row transformations. This is the precise separation
requested in the source, with a stronger two-sided obstruction as well.

\paragraph{Formal verification.}
The self-contained Lean 4.19.0 project defines actual integer matrices
as functions $\operatorname{Fin}2\to\operatorname{Fin}2\to\mathbb Z$.
Matrix multiplication, determinant, permanent, and unimodularity are
defined explicitly. Row equivalence and two-sided integer equivalence
quantify over \emph{all} integer matrices $U,V$; these quantifiers are
not replaced by a finite search or a bound on their entries.

The key invariant is the predicate that every matrix entry is an integer
multiple of two. The file proves that multiplication on either side by
any integer matrix preserves this predicate. It then verifies both
determinants, both permanents, the actual Smith-diagonal conditions, and
the inequivalences. The final existence theorem is
\texttt{conjecture\_6330} in \texttt{lean/Main.lean}.
It uses only bundled Std, with no \texttt{sorry}, \texttt{admit},
\texttt{native\_decide}, or added axiom. The printed axiom list is limited
to Lean's standard \texttt{propext} and \texttt{Quot.sound}.

\end{document}
